
\documentclass[11pt]{article}
\usepackage[margin=1in]{geometry}
\usepackage[T1]{fontenc}
\usepackage[utf8]{inputenc}
\usepackage{amsmath,amssymb,bm}
\usepackage{array,longtable,booktabs,graphicx}
\usepackage{enumitem}
\usepackage{hyperref}
\setlist[itemize]{leftmargin=1.5em}
\setlength{\parskip}{0.5em}
\setlength{\parindent}{0pt}
\renewcommand{\arraystretch}{1.2}

\begin{document}

\begin{center}
{\LARGE Combined LaTeX Transcriptions}\\[0.5em]
{\large Binder1.pdf Pages 1--20, followed by the first 9 standalone images}
\end{center}

This document combines the earlier LaTeX transcriptions into one editable \texttt{.tex} file.

\section*{Binder1.pdf --- Pages 1--20}

\subsection*{Page 1}
\section*{Planetary Rossby Waves}

Begin with the assumptions for conserving potential vorticity: inviscid, thin, ``rapid'' rotation.

\[
\frac{D}{Dt}\left(\frac{f+\zeta}{h}\right)=0
\]

The equations describing horizontal motions are:

\begin{align}
\frac{\partial u}{\partial t}
+u\frac{\partial u}{\partial x}
+v\frac{\partial u}{\partial y}
+w\frac{\partial u}{\partial z}
-fv
&=
-\frac{1}{\rho}\frac{\partial p}{\partial x}
\tag{1}\\[6pt]
\frac{\partial v}{\partial t}
+u\frac{\partial v}{\partial x}
+v\frac{\partial v}{\partial y}
+w\frac{\partial v}{\partial z}
+fu
&=
-\frac{1}{\rho}\frac{\partial p}{\partial y}
\tag{2}\\[6pt]
\frac{\partial u}{\partial x}
+\frac{\partial v}{\partial y}
+\frac{\partial w}{\partial z}
&=0
\tag{3}
\end{align}

Consider fluctuations with no mean flow (e.g.\ waves) and small velocities
(small $Ro$), so neglect nonlinear advection terms.

\subsection*{Page 2}
\begin{align}
\frac{\partial u}{\partial t}-fv
&=
-\frac{1}{\rho}\frac{\partial p}{\partial x}
\tag{1}\\[6pt]
\frac{\partial v}{\partial t}+fu
&=
-\frac{1}{\rho}\frac{\partial p}{\partial y}
\tag{2}
\end{align}

Substitute for $f=f_0+\beta y$,

\begin{align}
\frac{\partial u}{\partial t}-(f_0+\beta y)v
&=
-\frac{1}{\rho}\frac{\partial p}{\partial x}
\tag{1}\\[6pt]
\frac{\partial v}{\partial t}+(f_0+\beta y)u
&=
-\frac{1}{\rho}\frac{\partial p}{\partial y}
\tag{2}
\end{align}

\subsection*{Page 3}
\begin{align}
\frac{\partial u}{\partial t}-(f_0+\beta y)v
&=
-\frac{1}{\rho}\frac{\partial p}{\partial x}
\tag{1}\\[6pt]
\frac{\partial v}{\partial t}+(f_0+\beta y)u
&=
-\frac{1}{\rho}\frac{\partial p}{\partial y}
\tag{2}
\end{align}

Replace pressure gradients with surface height anomalies $\eta$,
that is,

\[
\frac{1}{\rho}\frac{\partial p}{\partial x}
=
g\frac{\partial \eta}{\partial x}
\qquad
\text{(hydrostatic)}
\]

so that

\begin{align}
\frac{\partial u}{\partial t}-(f_0+\beta y)v
&=
-g\frac{\partial \eta}{\partial x}
\tag{1}\\[6pt]
\frac{\partial v}{\partial t}+(f_0+\beta y)u
&=
-g\frac{\partial \eta}{\partial y}
\tag{2}
\end{align}

\subsection*{Page 4}
\begin{align}
\frac{\partial u}{\partial t}-(f_0+\beta y)v
&=
-g\frac{\partial \eta}{\partial x}
\tag{1}\\[6pt]
\frac{\partial v}{\partial t}+(f_0+\beta y)u
&=
-g\frac{\partial \eta}{\partial y}
\tag{2}
\end{align}

Recall that the flow is predominantly geostrophic, so the time-dependent
and beta terms are small compared to geostrophy.

As the geostrophic terms dominate, use these to express the velocities as

\[
u \approx -\frac{g}{f_0}\frac{\partial \eta}{\partial y},
\qquad
v \approx \frac{g}{f_0}\frac{\partial \eta}{\partial x}
\]

and substitute back into (1,2) for the small terms:

\begin{align}
-\frac{g}{f_0}\frac{\partial^2 \eta}{\partial y\,\partial t}
-f_0v
-\beta y \frac{g}{f_0}\frac{\partial \eta}{\partial x}
&=
-g\frac{\partial \eta}{\partial x}
\tag{4}\\[6pt]
\frac{g}{f_0}\frac{\partial^2 \eta}{\partial x\,\partial t}
+f_0u
-\beta y \frac{g}{f_0}\frac{\partial \eta}{\partial y}
&=
-g\frac{\partial \eta}{\partial y}
\tag{5}
\end{align}

\subsection*{Page 5}
Rearrange (4,5) for $v$ and $u$:

\begin{align}
v
&=
\frac{g}{f_0}\frac{\partial \eta}{\partial x}
-\frac{g}{f_0^2}\frac{\partial^2 \eta}{\partial y\,\partial t}
-\beta y \frac{g}{f_0^2}\frac{\partial \eta}{\partial x}
\tag{6}\\[6pt]
u
&=
-\frac{g}{f_0}\frac{\partial \eta}{\partial y}
-\frac{g}{f_0^2}\frac{\partial^2 \eta}{\partial x\,\partial t}
+\beta y \frac{g}{f_0^2}\frac{\partial \eta}{\partial y}
\tag{7}
\end{align}

The blue terms are ageostrophic. These describe the behaviour of the wave
perturbations superimposed on the geostrophic flow.

\subsection*{Page 6}
Recall (3), the continuity equation,

\[
\frac{\partial u}{\partial x}
+\frac{\partial v}{\partial y}
+\frac{\partial w}{\partial z}
=
0
\tag{3}
\]

This can be integrated over a layer depth $H$,

\[
\int_0^H
\left(
\frac{\partial u}{\partial x}
+\frac{\partial v}{\partial y}
+\frac{\partial w}{\partial z}
\right)\,dz
=
H\left(
\frac{\partial u}{\partial x}
+\frac{\partial v}{\partial y}
\right)
+w(H)-w(0)
=
0
\]

\[
\frac{\partial \eta}{\partial t}
+
H\left(
\frac{\partial u}{\partial x}
+\frac{\partial v}{\partial y}
\right)
=
0
\tag{8}
\]

\subsection*{Page 7}
Substitute $u,v$ from (6,7) into (8):

\begin{align}
v
&=
\frac{g}{f_0}\frac{\partial \eta}{\partial x}
-\frac{g}{f_0^2}\frac{\partial^2 \eta}{\partial y\,\partial t}
-\beta y \frac{g}{f_0^2}\frac{\partial \eta}{\partial x},
\\[6pt]
u
&=
-\frac{g}{f_0}\frac{\partial \eta}{\partial y}
-\frac{g}{f_0^2}\frac{\partial^2 \eta}{\partial x\,\partial t}
+\beta y \frac{g}{f_0^2}\frac{\partial \eta}{\partial y}
\tag{6,7}
\end{align}

\[
\frac{\partial \eta}{\partial t}
+
H\left(
\frac{\partial u}{\partial x}
+\frac{\partial v}{\partial y}
\right)
=
0
\tag{8}
\]

\begin{align*}
\frac{\partial \eta}{\partial t}
+\frac{gH}{f_0^2}
\Biggl(
\frac{\partial}{\partial x}
\left[
-f_0\frac{\partial \eta}{\partial y}
-\frac{\partial^2 \eta}{\partial x\,\partial t}
+\beta y \frac{\partial \eta}{\partial y}
\right]
+
\frac{\partial}{\partial y}
\left[
f_0\frac{\partial \eta}{\partial x}
-\frac{\partial^2 \eta}{\partial y\,\partial t}
-\beta y \frac{\partial \eta}{\partial x}
\right]
\Biggr)
&=0
\\[8pt]
\frac{\partial \eta}{\partial t}
+\frac{gH}{f_0^2}
\left(
-\frac{\partial^3 \eta}{\partial x^2\,\partial t}
-\frac{\partial^3 \eta}{\partial y^2\,\partial t}
-\beta \frac{\partial \eta}{\partial x}
\right)
&=0
\\[8pt]
\frac{\partial \eta}{\partial t}
+\frac{gH}{f_0^2}
\left(
-\frac{\partial}{\partial t}
\left[
\frac{\partial^2 \eta}{\partial x^2}
+\frac{\partial^2 \eta}{\partial y^2}
\right]
-\beta \frac{\partial \eta}{\partial x}
\right)
&=0
\end{align*}

\subsection*{Page 8}
\[
\frac{\partial \eta}{\partial t}
+\frac{gH}{f_0^2}
\left(
-\frac{\partial}{\partial t}
\left[
\frac{\partial^2 \eta}{\partial x^2}
+\frac{\partial^2 \eta}{\partial y^2}
\right]
-\beta \frac{\partial \eta}{\partial x}
\right)
=0
\]

Substitute for

\[
\nabla^2 \eta
=
\frac{\partial^2 \eta}{\partial x^2}
+
\frac{\partial^2 \eta}{\partial y^2}
\]

and define

\[
R=\frac{\sqrt{gH}}{f_0}
\]

as the Rossby radius:

\[
\frac{\partial \eta}{\partial t}
-
R^2\frac{\partial}{\partial t}\nabla^2\eta
-
R^2\beta\frac{\partial \eta}{\partial x}
=
0
\tag{9}
\]

This is the Rossby wave equation.

\subsection*{Page 9}
\[
\int_0^H
\left(
\frac{\partial u}{\partial x}
+\frac{\partial v}{\partial y}
+\frac{\partial w}{\partial z}
\right)\,dz
=
\frac{\partial \eta}{\partial t}
+
H\left(
\frac{\partial u}{\partial x}
+\frac{\partial v}{\partial y}
\right)
=
0
\]

\[
\frac{1}{\rho}\frac{\partial p}{\partial x}
=
g\frac{\partial \eta}{\partial x}
\qquad
\text{(hydrostatic)}
\]

\[
R=\frac{\sqrt{gH}}{f_0}
\qquad
\text{as the Rossby radius}
\]

\[
\frac{\partial \eta}{\partial t}
-
R^2\frac{\partial}{\partial t}\nabla^2\eta
-
R^2\beta\frac{\partial \eta}{\partial x}
=
0
\]

\subsection*{Page 10}
\[
\frac{\partial \eta}{\partial t}
-
R^2\frac{\partial}{\partial t}\nabla^2\eta
-
R^2\beta\frac{\partial \eta}{\partial x}
=
0
\tag{9}
\]

This is the Rossby wave equation.

The coefficients of (9) are constant, and we have assumed the perturbation
is a wave, so we prescribe the form

\[
\eta=\eta_0\cos(kx-my-\omega t).
\]

This allows us to rewrite (9) in terms of $\omega$ as a dispersion relation
for Rossby waves:

\[
\omega
=
-\beta R^2
\frac{k}{1+R^2(k^2+m^2)}
\]

\subsection*{Page 11}
Dispersion relation for planetary Rossby waves
\[
\left(
\text{recall } R=\sqrt{gH}/f_0
\right):
\]

\[
\omega
=
-\beta R^2
\frac{k}{1+R^2(k^2+m^2)}
\]

\begin{itemize}
\item For $\beta=0 \to \omega=0$: steady, geostrophic flow at constant $f$
(no waves).
\item We have assumed $Ro \ll 1$, so solution only valid for $\omega \ll f_0$.
\item Consider $k,m \sim 1/L \to \omega \sim \beta R^2 L/(L^2+R^2)$.
\item So for short waves $L<R$: $\omega \sim \beta L$.
\item For long waves $L>R$: $\omega \sim \beta R^2/L < \beta L$.
\item For $L=R$: $\omega$ is maximum at $\omega_{\max}=\beta R/2$.
\end{itemize}

\subsection*{Page 12}
Dispersion relation for planetary Rossby waves
\[
\left(
\text{recall } R=\sqrt{gH}/f_0
\right):
\]

\[
\omega
=
-\beta R^2
\frac{k}{1+R^2(k^2+m^2)}
\]

\begin{itemize}
\item ZONAL phase speed always negative $\to$ westward propagation
\[
c_x=\frac{\omega}{k}
=
-\beta R^2
\frac{1}{1+R^2(k^2+m^2)}
\]
\item MERIDIONAL phase speed either sign
\[
c_y=\frac{\omega}{m}
=
-\beta R^2
\frac{k/m}{1+R^2(k^2+m^2)}
\]
\item Very long waves $(1/k,1/m,L \gg R)$: propagate westward at
\[
c=-\beta R^2,
\]
maximum speed allowed.
\end{itemize}

\subsection*{Page 13}
\[
\frac{D}{Dt}\left(\frac{f+\zeta}{h}\right)=0
\]

\[
\text{North, high }f
\qquad\qquad
\text{South, low }f
\]

\[
\text{Contours of }f
\]

\[
\text{Greater }f,\ \text{reduced }\zeta
\qquad\qquad
\text{Smaller }f,\ \text{increased }\zeta
\]

\[
\text{Greater }f,\ \text{reduced }\zeta
\]

\subsection*{Page 14}
\section*{Taylor--Proudman Theorem}

For steady, $Ro \ll 1$ and uniform density $\rho$.

Geostrophic vector equation:

\[
2\Omega \times \mathbf{u}
=
-\frac{\nabla p}{\rho}
\]

Take the curl of this equation:

\[
\nabla \times (2\Omega \times \mathbf{u})
=
-\nabla \times \left(\frac{\nabla p}{\rho}\right)
\]

RHS: curl of a gradient $(\nabla p)$ is zero, $\rho$ is uniform.

So:

\[
\nabla \times (\Omega \times \mathbf{u})=0
\]

Expands to:

\[
(\Omega\cdot\nabla)\mathbf{u}
-
(\mathbf{u}\cdot\nabla)\Omega
+
\mathbf{u}(\nabla\cdot\Omega)
-
\Omega(\nabla\cdot\mathbf{u})
=
0
\]

\subsection*{Page 15}
\section*{Taylor--Proudman Theorem}

\[
(\Omega\cdot\nabla)\mathbf{u}
-
(\mathbf{u}\cdot\nabla)\Omega
+
\mathbf{u}(\nabla\cdot\Omega)
-
\Omega(\nabla\cdot\mathbf{u})
=
0
\]

$\Omega$ is uniform everywhere, and continuity is unaltered by rotation, so:

\[
(\Omega\cdot\nabla)\mathbf{u}=0
\]

Choose coordinates such that $\Omega$ is in the $z$-direction,

\[
\Omega \frac{\partial \mathbf{u}}{\partial z}=0
\qquad\to\qquad
\frac{\partial u}{\partial z}
=
\frac{\partial v}{\partial z}
=
\frac{\partial w}{\partial z}
=
0
\]

And for a domain with a solid (bottom) boundary where $w=0$,

\[
\frac{\partial u}{\partial z}
=
\frac{\partial v}{\partial z}
=
0,
\qquad
w=0
\]

Horizontal velocities constant with depth, moves as ``Taylor columns''.

\subsection*{Page 16}
\section*{Flow on a Rotating Planet -- Horizontal Scaling}

Euler Equation in Rotating Frame:

\[
\frac{D\mathbf{u}}{Dt}+2\Omega\times\mathbf{u}
=
\mathbf{g}-\frac{\nabla p}{\rho}
\]

$x$-component of the Euler equation:

\[
\frac{\partial u}{\partial t}
+u\frac{\partial u}{\partial x}
+v\frac{\partial u}{\partial y}
+w\frac{\partial u}{\partial z}
+2\Omega w\cos\theta
-2\Omega v\sin\theta
=
-\frac{1}{\rho}\frac{\partial p}{\partial x}
\]

Assume hydrostatic pressure scale:

\[
\Delta P \sim \rho' g H
\]

horizontal velocity scale $U$, timescale $T=L/U$,
\[
w \sim \frac{UH}{L}
\]

Scales of terms:

\[
\frac{U^2}{L}
+\frac{U^2}{L}
+\frac{U^2}{L}
+\frac{U^2}{L}
+2\Omega\frac{UH}{L}
-2\Omega U
\sim
\frac{\rho' g H}{\rho L}
\]

\subsection*{Page 17}
\section*{Flow on a Rotating Planet -- Horizontal Scaling in $x$}

$x$-component of the Euler equation:

\[
\frac{\partial u}{\partial t}
+u\frac{\partial u}{\partial x}
+v\frac{\partial u}{\partial y}
+w\frac{\partial u}{\partial z}
+2\Omega w\cos\theta
-2\Omega v\sin\theta
=
-\frac{1}{\rho}\frac{\partial p}{\partial x}
\]

\[
\frac{U^2}{L}
+\frac{U^2}{L}
+\frac{U^2}{L}
+\frac{U^2}{L}
+2\Omega\frac{UH}{L}
-2\Omega U
\sim
\frac{\rho' g H}{\rho L}
\]

\[
1+1+1+1
+
\frac{2\Omega L}{U}\left(\frac{H}{L}-1\right)
\sim
\frac{\rho' g H}{\rho U^2}
\]

Rossby Number:
\[
Ro=\frac{U}{2\Omega L}
\]

\[
\to \text{ ratio of inertial to rotation timescales}
\]

\[
Ro \ll 1
\quad
\text{for ``slow'' flow, ``large'' lengthscales and ``rapid'' rotation rates}
\]

\[
Ro \ll 1
\qquad\to\qquad
\text{Coriolis balances pressure gradient}
\]

\[
2\Omega v\sin\theta
=
\frac{1}{\rho}\frac{\partial p}{\partial x}
\]

\subsection*{Page 18}
\section*{Flow on a Rotating Planet -- Horizontal Scaling in $y$}

$y$-component of the Euler equation:

\[
\frac{\partial v}{\partial t}
+u\frac{\partial v}{\partial x}
+v\frac{\partial v}{\partial y}
+w\frac{\partial v}{\partial z}
+2\Omega u\sin\theta
=
-\frac{1}{\rho}\frac{\partial p}{\partial y}
\]

Scales of terms:

\[
\frac{U^2}{L}
+\frac{U^2}{L}
+\frac{U^2}{L}
+\frac{U^2}{L}
+2\Omega U
\sim
\frac{\rho' g H}{\rho L}
\]

\[
1+1+1+1+\frac{2\Omega L}{U}
\sim
\frac{\rho' g H}{\rho U^2}
\]

Again, Rossby Number:
\[
Ro=\frac{U}{2\Omega L}
\]

\[
\to \text{ ratio of inertial to rotation timescales}
\]

\[
Ro \ll 1
\quad
\text{for ``slow'' flow, ``large'' length scales and ``rapid'' rotation rates}
\]

\[
Ro \ll 1
\qquad\to\qquad
\text{Coriolis balances pressure gradient}
\]

\[
-2\Omega u\sin\theta
=
\frac{1}{\rho}\frac{\partial p}{\partial y}
\]

\subsection*{Page 19}
\section*{Navier--Stokes Momentum Equation in Rotating Frame}

\[
\frac{D\mathbf{u}}{Dt}
+2\Omega\times\mathbf{u}
=
\mathbf{g}
-\frac{\nabla p}{\rho}
+\nu\nabla^2\mathbf{u}
\]

\section*{Euler Equation in Rotating Frame}

\[
\frac{D\mathbf{u}}{Dt}
+2\Omega\times\mathbf{u}
=
\mathbf{g}
-\frac{\nabla p}{\rho}
\]

\[
2\Omega\times\mathbf{u}
=
(-fv,\,fu,\,0)
\]

\[
f=2\Omega\sin\theta
\]

\subsection*{Page 20}
\section*{Flow on a Rotating Planet -- Vertical Scaling}

Euler Equation in Rotating Frame:

\[
\frac{D\mathbf{u}}{Dt}+2\Omega\times\mathbf{u}
=
\mathbf{g}-\frac{\nabla p}{\rho}
\]

Vertical component of the Euler equation:

\[
\frac{\partial w}{\partial t}
+u\frac{\partial w}{\partial x}
+v\frac{\partial w}{\partial y}
+w\frac{\partial w}{\partial z}
-2\Omega u\cos\theta
=
-g-\frac{1}{\rho}\frac{\partial p}{\partial z}
\]

\[
H \ll L
\]

\[
T=\frac{L}{U}
\]

\[
\frac{U}{L}\sim\frac{W}{H}
\qquad\Rightarrow\qquad
W\sim \frac{UH}{L}
\]

\clearpage
\section*{Standalone Images 1--9}

\subsection*{Image 1}
This is the $1^{\text{st}}$ generation only. We have two more copies of all these rows, identical in every way except for mass. $2^{\text{nd}}$ and $3^{\text{rd}}$ gen matter is unstable!

\begin{center}
\small
\resizebox{\textwidth}{!}{%
\begin{tabular}{|c|c|c|}
\hline
\textbf{$1^{\text{st}}$ Gen} & \textbf{$2^{\text{nd}}$ Gen} & \textbf{$3^{\text{rd}}$ Gen} \\
\hline
Electron $e$ & Muon $\mu$ & Tau $\tau$ \\
\hline
Electron Neutrino $\nu_e$ & Muon Neutrino $\nu_\mu$ & Tau Neutrino $\nu_\tau$ \\
\hline
Up Quark $u$ & Charm Quark $c$ & Top (Truth) Quark $t$ \\
\hline
Down Quark $d$ & Strange Quark $s$ & Bottom (Beauty) Quark $b$ \\
\hline
\end{tabular}%
}
\end{center}

\[
T_3 \equiv \tau_3/2
\]

\[
\text{So, visually it's like:}
\]

\[
Q = \frac{1}{2}\left(\tau_3 + Y\right)
\]

\[
\text{rotate around z-axis}
\qquad
\text{rotate around flagpole}
\]

\begin{center}
\small
\resizebox{\textwidth}{!}{%
\begin{tabular}{|l|c|c|c|}
\hline
\textbf{Particle} & \textbf{Electric Charge $Q$} & \textbf{Weak Isospin $T_3$} & \textbf{Hypercharge $Y$} \\
\hline
Electron $e_R$ & $-1$ & $0$ & $-2$ \\
Neutrino $\nu_e$ & $0$ & $+1/2$ & $-1$ \\
Electron $e_L$ & $-1$ & $-1/2$ & $-1$ \\
\hline
Up Quark $u_R$ & $+2/3$ & $0$ & $+4/3$ \\
Down Quark $d_R$ & $-1/3$ & $0$ & $-2/3$ \\
Up Quark $u_L$ & $+2/3$ & $+1/2$ & $+1/3$ \\
Down Quark $d_L$ & $-1/3$ & $-1/2$ & $+1/3$ \\
\hline
$W^\pm$ Bosons & $\pm 1$ & $\pm 1$ & $0$ \\
$Z$ Boson & $0$ & $0$ & $0$ \\
Photon $\gamma$ & $0$ & $0$ & $0$ \\
Gluon & $0$ & $0$ & $0$ \\
\hline
Higgs $H^+$ & $+1$ & $+1/2$ & $+1$ \\
Higgs $H^0$ & $0$ & $-1/2$ & $+1$ \\
\hline
\end{tabular}%
}
\end{center}

\[
\text{Leptons}
\qquad
\text{Quarks}
\qquad
\text{Gauge Bosons}
\qquad
\text{Higgs Field}
\]

\[
\text{$SU(2)_L$ doublets: }
\begin{pmatrix}
\nu_e \\
e_L
\end{pmatrix},
\quad
\begin{pmatrix}
u_L \\
d_L
\end{pmatrix},
\quad
\begin{pmatrix}
H^+ \\
H^0
\end{pmatrix}
\]

\[
\text{Gell-Mann--Nishijima relation:}
\qquad
Q = T_3 + \frac{Y}{2}
\]

\subsection*{Image 2}
\[
\Psi =
\begin{pmatrix}
\psi_r \\
\psi_g \\
\psi_b
\end{pmatrix},
\qquad
\bar{\Psi} =
\begin{pmatrix}
\bar{\psi}_r & \bar{\psi}_g & \bar{\psi}_b
\end{pmatrix}
\]

\subsection*{Image 3}
\[
\text{The color triplet field } \Psi \in \mathbb{C}^3 \otimes \mathbb{C}^4 \cong \mathbb{C}^{12}
\text{ gives us a way to write } \mathcal{L} \text{ more concisely, but now that}
\]

\[
\text{we're here... what happens if we require } \mathcal{L}
\text{ to remain unchanged under any transformation}
\]

\[
\Psi \to \Psi' = U(x)\Psi
\qquad
\text{where } U(x) \in SU(3)\text{?}
\]

\subsection*{Image 4}
\[
\text{Free quark, non-interacting three-color Dirac lagrangian:}
\]

\begin{align*}
\mathcal{L} &= \mathcal{L}_r + \mathcal{L}_g + \mathcal{L}_b \\
\mathcal{L}
&= \bar{\psi}_r \left(i\gamma^\mu \partial_\mu - m\right)\psi_r
 + \bar{\psi}_g \left(i\gamma^\mu \partial_\mu - m\right)\psi_g
 + \bar{\psi}_b \left(i\gamma^\mu \partial_\mu - m\right)\psi_b
\end{align*}

\[
\mathcal{L} = \bar{\Psi}\left(i\gamma^\mu \partial_\mu - m\right)\Psi
\]

\subsection*{Image 5}
\[
\text{Free quark, non-interacting three-color Dirac lagrangian:}
\]

\begin{align*}
\mathcal{L} &= \mathcal{L}_r + \mathcal{L}_g + \mathcal{L}_b \\
\mathcal{L}
&= \bar{\psi}_r \left(i\gamma^\mu \partial_\mu - m\right)\psi_r \\
&\quad + \bar{\psi}_g \left(i\gamma^\mu \partial_\mu - m\right)\psi_g \\
&\quad + \bar{\psi}_b \left(i\gamma^\mu \partial_\mu - m\right)\psi_b
\end{align*}

\subsection*{Image 6}
\begin{align*}
\psi_r &= u_s(\vec{p}) \exp\!\left[-i\left(Et - \vec{p}\cdot\vec{x}\right)\right], \\
\psi_g &= u_s(\vec{p}) \exp\!\left[-i\left(Et - \vec{p}\cdot\vec{x}\right)\right], \\
\psi_b &= u_s(\vec{p}) \exp\!\left[-i\left(Et - \vec{p}\cdot\vec{x}\right)\right].
\end{align*}

\subsection*{Image 7}
\begin{center}
\small
\resizebox{\textwidth}{!}{%
\begin{tabular}{|c|c|c|c|c|c|c|}
\hline
\textbf{Baryon} &
\textbf{Quark content} &
\textbf{Charge} &
\shortstack{\textbf{Mass}\\\textbf{$(\mathrm{GeV}/c^2)$}} &
\shortstack{\textbf{Half Life}\\\textbf{(seconds)}} &
\textbf{Principal decays} &
\textbf{Spin} \\
\hline
$p$ & $uud$ & $+1$ & $0.93828$ & at least $3\cdot10^{41}$ & -- & $1/2$ \\
$n$ & $udd$ & $0$ & $0.93957$ & $880$ & $pe\bar{\nu}_e$ & $1/2$ \\
$\Lambda$ & $uds$ & $0$ & $1.1156$ & $2.63\cdot10^{-10}$ & $p\pi^{-},\, n\pi^{0}$ & $1/2$ \\
$\Sigma^{+}$ & $uus$ & $+1$ & $1.1894$ & $0.80\cdot10^{-10}$ & $p\pi^{0},\, n\pi^{+}$ & $1/2$ \\
$\Sigma^{0}$ & $uds$ & $0$ & $1.1925$ & $6\cdot10^{-20}$ & $\Lambda\gamma$ & $1/2$ \\
$\Sigma^{-}$ & $dds$ & $-1$ & $1.1973$ & $1.48\cdot10^{-10}$ & $n\pi^{-}$ & $1/2$ \\
$\Xi^{0}$ & $uss$ & $0$ & $1.3149$ & $2.90\cdot10^{-10}$ & $\Lambda\pi^{0}$ & $1/2$ \\
$\Xi^{-}$ & $dss$ & $-1$ & $1.3213$ & $1.64\cdot10^{-10}$ & $\Lambda\pi^{-}$ & $1/2$ \\
$\Lambda_c^{+}$ & $udc$ & $+1$ & $2.281$ & $2\cdot10^{-13}$ & $pK^{-}\pi^{+}$ & $1/2$ \\
$\Delta$ & $uuu,\, uud,\, udd,\, ddd$ & $+2,\,+1,\,0,\,-1$ & $1.232$ & $0.6\cdot10^{-23}$ & $N\pi$ & $3/2$ \\
$\Sigma^{*}$ & $uus,\, uds,\, dds$ & $+1,\,0,\,-1$ & $1.385$ & $2\cdot10^{-23}$ & $\Lambda\pi,\, \Sigma\pi$ & $3/2$ \\
$\Xi^{*}$ & $uss,\, dss$ & $0,\,-1$ & $1.533$ & $7\cdot10^{-23}$ & $\Xi\pi$ & $3/2$ \\
$\Omega^{-}$ & $sss$ & $-1$ & $1.672$ & $0.82\cdot10^{-10}$ & $\Lambda K^{-},\, \Xi^{0}\pi^{-},\, \Xi^{-}\pi^{0}$ & $3/2$ \\
\hline
\end{tabular}%
}
\end{center}

\textbf{Source:} \textit{Introduction to Elementary Particles}, David Griffiths, table at the beginning of the book.

\subsection*{Image 8}
\begin{center}
\small
\resizebox{\textwidth}{!}{%
\begin{tabular}{|c|c|c|c|c|c|}
\hline
\shortstack{\textbf{Vector}\\\textbf{Meson}\\\textbf{(Spin-1)}} &
\textbf{Quark content} &
\textbf{Charge} &
\shortstack{\textbf{Mass}\\\textbf{$(\mathrm{GeV}/c^2)$}} &
\shortstack{\textbf{Half Life}\\\textbf{(seconds)}} &
\textbf{Principal decays} \\
\hline
$\rho$ & $u\bar{d},\, d\bar{u},\, (u\bar{u}-d\bar{d})/\sqrt{2}$ & $+1,\,-1,\,0$ & $0.770$ & $0.4\cdot10^{-23}$ & $\pi\pi$ \\
$K^{*}$ & $u\bar{s},\, s\bar{u},\, d\bar{s},\, s\bar{d}$ & $+1,\,-1,\,0,\,0$ & $0.892$ & $1\cdot10^{-23}$ & $K\pi$ \\
$\omega$ & $(u\bar{u}+d\bar{d})/\sqrt{2}$ & $0$ & $0.783$ & $7\cdot10^{-23}$ & $\pi^{+}\pi^{-}\pi^{0},\, \pi^{0}\gamma$ \\
$\phi$ & $s\bar{s}$ & $0$ & $1.020$ & $20\cdot10^{-23}$ & $K^{+}K^{-},\, K^{0}\bar{K}^{0}$ \\
$J/\psi$ & $c\bar{c}$ & $0$ & $3.097$ & $1\cdot10^{-20}$ & $e^{+}e^{-},\, \mu^{+}\mu^{-},\, 5\pi,\, 7\pi$ \\
$D^{*}$ & $c\bar{d},\, d\bar{c},\, c\bar{u},\, u\bar{c}$ & $+1,\,-1,\,0,\,0$ & $2.010$ & $>1\cdot10^{-22}$ & $D\pi,\, D\gamma$ \\
$\Upsilon$ & $b\bar{b}$ & $0$ & $9.460$ & $2\cdot10^{-20}$ & $\tau^{+}\tau^{-},\, \mu^{+}\mu^{-},\, e^{+}e^{-}$ \\
\hline
\end{tabular}%
}
\end{center}

\subsection*{Image 9}
\begin{center}
\small
\resizebox{\textwidth}{!}{%
\begin{tabular}{|c|c|c|c|c|c|}
\hline
\shortstack{\textbf{Pseudoscalar}\\\textbf{Meson}\\\textbf{(Spin-0)}} &
\textbf{Quark content} &
\textbf{Charge} &
\shortstack{\textbf{Mass}\\\textbf{$(\mathrm{GeV}/c^2)$}} &
\shortstack{\textbf{Half Life}\\\textbf{(seconds)}} &
\textbf{Principal decays} \\
\hline
$\pi^{\pm}$ & $u\bar{d},\, d\bar{u}$ & $+1,\,-1$ & $0.139569$ & $2.60\cdot10^{-8}$ & $\mu\nu_{\mu}$ \\
$\pi^{0}$ & $(u\bar{u}-d\bar{d})/\sqrt{2}$ & $0$ & $0.134964$ & $8.7\cdot10^{-17}$ & $\gamma\gamma$ \\
$K^{\pm}$ & $u\bar{s},\, s\bar{u}$ & $+1,\,-1$ & $0.49367$ & $1.24\cdot10^{-8}$ & $\mu\nu_{\mu},\, \pi^{\pm}\pi^{0},\, \pi^{\pm}\pi^{\pm}\pi^{\mp}$ \\
$K^{0},\, \bar{K}^{0}$ & $d\bar{s},\, s\bar{d}$ & $0,\,0$ & $0.49772$ & $K_S^{0}: 0.892\cdot10^{-10}$ & $\pi^{+}\pi^{-},\, \pi^{0}\pi^{0}$ \\
$K^{0},\, \bar{K}^{0}$ & $d\bar{s},\, s\bar{d}$ & $0,\,0$ & $0.49772$ & $K_L^{0}: 5.18\cdot10^{-8}$ & $\pi e \nu_e,\, \pi\mu\nu_{\mu},\, \pi\pi\pi$ \\
$\eta$ & $(u\bar{u}+d\bar{d}-2s\bar{s})/\sqrt{6}$ & $0$ & $0.5488$ & $7\cdot10^{-19}$ & $\gamma\gamma,\, \pi^{0}\pi^{0}\pi^{0},\, \pi^{+}\pi^{-}\pi^{0}$ \\
$\eta'$ & $(u\bar{u}+d\bar{d}+s\bar{s})/\sqrt{3}$ & $0$ & $0.9576$ & $3\cdot10^{-21}$ & $\eta\pi\pi,\, \rho^{0}\gamma$ \\
$D^{\pm}$ & $c\bar{d},\, d\bar{c}$ & $+1,\,-1$ & $1.869$ & $9\cdot10^{-13}$ & $K\pi\pi$ \\
$D^{0},\, \bar{D}^{0}$ & $c\bar{u},\, u\bar{c}$ & $0,\,0$ & $1.865$ & $4\cdot10^{-13}$ & $K\pi\pi$ \\
$D_s^{\pm}$ & $c\bar{s},\, s\bar{c}$ & $+1,\,-1$ & $1.971$ & $3\cdot10^{-13}$ & \text{not established} \\
$B^{\pm}$ & $u\bar{b},\, b\bar{u}$ & $+1,\,-1$ & $5.271$ & $14\cdot10^{-13}$ & $D+?$ \\
$B^{0},\, \bar{B}^{0}$ & $d\bar{b},\, b\bar{d}$ & $0,\,0$ & $5.275$ & $14\cdot10^{-13}$ & $D+?$ \\
$\eta_c$ & $c\bar{c}$ & $0$ & $2.981$ & $6\cdot10^{-23}$ & $KK\pi,\, \eta\pi\pi,\, \eta'\pi\pi$ \\
\hline
\end{tabular}%
}
\end{center}

\end{document}
